\section{Практическая работа №4}
\label{sec:pr4}

\subsection{Выполненные задания}

В четвёртой практической работе был создан Stack-проект \mintinline{text}{mazeGame}. Сразу после инициализации был удалён стандартный файл \mintinline{text}{Lib}. В проект добавлены зависимости, файл с лабиринтом главного здания, функция чтения файла в формат \mintinline{haskell}{IO [(String, String)]}, собственный трансформер \mintinline{haskell}{MyStateT} и функция запуска игры.

Лабиринт представлен типом \mintinline{haskell}{Map Room [Room]}, где ключ --- комната, а значение --- список соседних комнат. Функция \mintinline{haskell}{addRoomPair} добавляет связь между двумя комнатами в обе стороны, поэтому из списка пар строится неориентированный граф.

Тип игрового действия после рефакторинга:

\captionof{listing}{Стек трансформеров игры}
\label{lst:game-type-pr4}
\begin{minted}{haskell}
type Game = ReaderT Maze (WriterT [String] (MyStateT Room IO))
\end{minted}

Основные коммиты практики: \mintinline{text}{e4678a9}, \mintinline{text}{94e6e47}, \mintinline{text}{1f73730}, \mintinline{text}{bba7573}, \mintinline{text}{95f97e2}, \mintinline{text}{cbaa959}, \mintinline{text}{99277ac}, \mintinline{text}{515420e}.

\subsection{Использование нейросетей}

Для четвёртой практики был создан файл с перепиской с LLM. После завершения работы лог больше не дополнялся. Модель использовалась для обсуждения трансформера, порядка терминирования и объяснения функций \mintinline{haskell}{lift}, \mintinline{haskell}{ask}, \mintinline{haskell}{getS}, \mintinline{haskell}{runGame}.

\subsection{Вопросы защиты и их решения}

Трансформер \mintinline{haskell}{MyStateT} объявлен в отдельном файле \mintinline{text}{MyStateT.hs}. Он является конструктором, который принимает монаду и возвращает новую монаду с добавленным состоянием и сохранённым функционалом внутренней монады.

Терминирование трансформеров происходит в функции \mintinline{haskell}{runGame}. Каждый слой раскрывается своей функцией \mintinline{haskell}{runReaderT}, \mintinline{haskell}{runWriterT}, \mintinline{haskell}{runMyStateT}; после раскрытия всех слоёв остаётся терминатор \mintinline{haskell}{IO}.

\captionof{listing}{Запуск игры}
\label{lst:run-game-pr4}
\begin{minted}{haskell}
runMyStateT (runWriterT (runReaderT game (buildMaze roomPairs))) "start"
\end{minted}

В этом выражении \mintinline{haskell}{game} передаётся в \mintinline{haskell}{runReaderT}, туда же передаётся окружение --- \mintinline{haskell}{buildMaze roomPairs}. После этого раскрывается \mintinline{haskell}{WriterT}, затем \mintinline{haskell}{MyStateT} получает начальное состояние \mintinline{haskell}{"start"}. \mintinline{haskell}{IO} вызывается не самим написанием выражения, а при выполнении полученного \mintinline{haskell}{IO}-действия, например через bind или в \mintinline{haskell}{main}.

\newpage
